% Includable proof fragment. No changes to the assembled paper.
\subsection{A redundant denominator guard in every pole block}

\begin{lemma}\label{lem:pole-denominator-redundant}
Let $F$ be a field of characteristic different from $2$ and $3$,
and let $c\in F$ satisfy $-c\notin F^2$.
For $b,t,e,p,q,s\in F$, put
\[
 \delta=1+bt^2,\qquad B=cb,\qquad
 K=B\delta^2-b^2(1-bt^2)^2.
\]
If
\begin{align}
 p^2&=B(1-e)(1+3e),\label{eq:pole-first-cleared}\\
 Bq^2\delta^4
   &=(s^2\delta^2-BeK)(s^2\delta^2+3BeK),
       \label{eq:pole-second-cleared}
\end{align}
then $\delta\ne0$.
Consequently these two polynomial equations are equivalent to the
two rational pole equations obtained by setting
\[
 U=\frac{b(1-bt^2)}{1+bt^2},\qquad h=e(B-U^2)
\]
and requiring
\[
 p^2=B(1-e)(1+3e),\qquad
 Bq^2=(s^2-Bh)(s^2+3Bh).
\]
\end{lemma}
\begin{proof}
The nonsquare hypothesis implies $c\ne0$.
If $\delta=0$, then $b,t\ne0$, $bt^2=-1$ and $K=-4b^2$.
Equation~\eqref{eq:pole-second-cleared} becomes
\[
 0=-48 B^2b^4e^2.
\]
All its factors except possibly $e$ are nonzero, so $e=0$.
Equation~\eqref{eq:pole-first-cleared} now gives
$p^2=cb=-c/t^2$, hence $(pt)^2=-c$, a contradiction.
When $\delta\ne0$, the identity $h=eK/\delta^2$ shows that
\eqref{eq:pole-second-cleared} is precisely the second rational
equation multiplied by $\delta^4$.  This proves both directions.
\end{proof}

For the construction over $\mathbb Q$ one has $c=17$; the hypothesis
holds because $-17$ is not a rational square.  Thus no separate inverse
is needed for any of these $42$ denominators.

The argument is a statement about rational solutions.  It must not be
silently interpreted as an equivalence of geometric fibers over an
algebraic closure, where $-17$ becomes a square.

\paragraph{Compatibility with reduced circuit numerators.}
It remains to justify deleting the guards from the existing rational
arithmetic circuit, whose fractions were simplified before their
numerators were stored.  In a fixed pole block write
$b=N/Q$ and $\Delta=Q+Nt^2$, so that $\delta=\Delta/Q$.
The retained common guard ensures $Q\ne0$.
The exact circuit factor check gives the stored second pole residual
as $n/(S\Delta^4)$, where $S$ is a nonzero constant times a product
of retained guard factors.  On the retained open set, the rational
identity defining the circuit therefore gives
\[
 \delta^4\left(Bq^2-(s^2-Bh)(s^2+3Bh)\right)
       =\frac{n}{S Q^4}.
\]
The expression on the left is the polynomial expression
\eqref{eq:pole-second-cleared}, interpreted with $b=N/Q$.
Both sides are regular on the retained open set, including
$\Delta=0$; the equality follows by multiplying denominators and
using a polynomial identity.  Hence $n=0$ implies the cleared pole
equation even at a putative zero of the deleted guard.
The first residual has denominator $Q$, and consequently supplies
\eqref{eq:pole-first-cleared} there as well.
Lemma~\ref{lem:pole-denominator-redundant} excludes that zero.
This proves that cancellation of rational expressions before
numerator extraction creates no gap in the guard deletion.

\paragraph{Verified independence of the remaining guard.}
An independent traversal of the frozen $855$-variable circuit finds
exactly $42$ guarded numerators involving pole variables, one for
each pole coordinate $t$.  They are the $42$ numerators $\Delta$.
Deleting them leaves $1639$ recorded guards, with $1560$ distinct
nonconstant factors.  Their combined variable support consists of
the input, all $633$ point coordinates, the coordinate denominator,
and the two contact variables.  In particular, the remaining common
guard is independent of all three proxy variables, all three anchor
variables, $\theta$, and every pole variable.

For every second pole residual, its denominator has its own
$\Delta$ to exponent exactly four and otherwise only retained
factors.  The other $255$ rational residuals have denominators
involving only retained factors.  These are exact structural checks;
the generic clearing and exceptional-fiber identities are also
verified symbolically by
\texttt{verify\_pole\_guard\_removal.py}, with receipt
\texttt{pole-guard-removal-verification.json}.
